\documentclass{report}
\usepackage{amsfonts, amsmath}
\newcommand\R{{\mathbb R}}
\newcommand\Q{{\mathbb Q}}
\newcommand\N{{\mathbb N}}
\newcommand\ve{\varepsilon}
\newcommand\Co{{\mathcal C}}
\pagestyle{empty}
\begin{document}
\large
\centerline{\Huge Fisica Matematica II}
\renewcommand{\thefootnote}{ } 
\footnote{\sl Avete 2 ore di tempo. Ogni esercizio
vale dieci punti. La sufficienza si ottiene con un punteggio $\geq$
18. Solo le risposte chiaramente giustificate saranno prese in
considerazione.}
\renewcommand{\thefootnote}{\arabic{footnote}} 
\vskip.1cm
\centerline{\Large Appello 17-06-2003}
\vskip1cm
\begin{enumerate}
\item Si determini la soluzione di
\[
\begin{split}
&xuu_x+u_y=1\\
&u(1,y)=\frac y2
\end{split}
\]
\item  Si consideri la soluzione dell'equazione
\[
\begin{split}
&u_t-u_{xx}=f\quad x\in[0,\pi]\\
&u(0,t)=u(\pi,t)=0\quad t>0\\
&u(x,0)=h(x)\quad x\in[0,\pi].
\end{split}
\]
Sia $v(x)=\lim_{t\to\infty}u(x,t)$. Come dipende $v$ da $h$. Si scriva
una equazione differenziale di cui $v$ \`e soluzione.
\item  Dati due vettori $v,w\in\R^2$, linearmente indipendenti, si
classifichino tutte le funzioni armoniche $f:\R^2\to\R$ tali che, per
ogni $\xi\in\R^2$,
\[
f(\xi+v)=f(\xi)=f(\xi+w).
\]
\item Si risolva l'equazione
\[
\begin{split}
&u_{xx}+\sin x=u_{tt}\quad x>0,\;t>0\\
&u(0,t)=0\quad t>0\\
&u(x,0)=x;\; u_t(x,0)=0\quad x\geq 0 
\end{split}
\]
\end{enumerate}
\end{document}
